% !TeX root = RJwrapper.tex
\title{\pkg{sbtools}: A package connecting R to cloud-based data for collaborative online research}
\author{by Luke A Winslow, Scott Chamberlain, Alison Appling and Jordan Read}

\maketitle

\abstract{
The adoption of high-quality tools
for collaboration and reproducibile research such as R and Github is becoming 
more common in many research fields. 
While Github and other version management systems are excellent 
resources, they were originally designed to handle code and scale poorly to 
large text-based or binary datasets. While a number of scientific data repositories 
are coming online, they are often focused on dataset archival and publication. 
To handle collaborative workflows using large scientific datasets, there is increasing 
need to connect cloud-based online data storage to R. In this article, we 
describe the new R package \pkg{sbtools} enables direct access to the advanced online data 
functionality provided by ScienceBase, the U.S. Geological Survey's online scientific data storage
platform. 
}

\section{Introduction}

Cloud data storage platforms can be a powerful tool for research collaboration,
distributed computing, and data publication. However, while browser-based graphical
interfaces make these platforms accessible to a wide audience, there is also
a strong need for scripted data handling so that researchers can rigorously
document data provenance and create reproducible analyses. For example,
Figshare \cidep{figshare}, DataOne \cidep{dataone}, and Dataverse \cidep{dataverse} 
are all research data
sharing platforms gaining use in different research fields. For each of these
platforms, the community has released \citep{rfigshare, dvn} or is
developing \citep{dataone} R packages to streamline the storage and access
to archived data. 

These existing projects primarily focus on creating more useful, open, and
accessible the end products of research. But data- and code-intensive collaborative projects
increasingly need collaborative solutions for data storage, sharing, and updating
not just at the end of the project, but throughout the full project lifecycle. 
Github is increasingly used to collaborate around data products \citep{GandrudGithub}, 
but does not scale well to the distribution of large datasets \citep{Delcambre2013} and does not
include metadata or queries beyond free-text search. Furthermore, 
most data archive platforms like DataOne and Dataverse 
focus on academic research projects and do not support some
of the unique needs of federal research. For example, many government institutions
require an archived copy of released data to be stored and available through federally
operated websites. Third party storage providers are insufficient.

To address future scientific data sharing and archival challenges, the USGS
created ScienceBase (\url{https://www.sciencebase.gov/}). This platform is
designed to support
the full project data lifecycle and has seen rapid adoption with USGS
researchers and collaborators. ScienceBase supports the storage and access
of large files and datasets. It allows data to be stored with a user-configurable 
mixture of
public and controlled access and has been designed from the beginning with
first-class RESTful web interfaces to support robust API access. Items on ScienceBase 
can have hierarchical relationships, facilitating the organization of complex or related
data. It also allows
a seamless transition from project development to data publication, with
a focus on complete metadata description, search and public data access once 
data have been released.

To expand the usefulness of the ScienceBase platform directly to R workflows, 
we have designed and
implemented an R interface to ScienceBase called \pkg{sbtools}.
This interface provides scripted R access to ScienceBase to manage metadata and
data files, to search the catalog of datasets, and to view and modify data in
formats familiar to R users.
Here we describe unique and useful features of ScienceBase, and how we have
implemented the R interface to make them accessible and useful from R.


\section{USGS ScienceBase}

To facilitate and encourage data sharing and dissemination, 
the U.S. Geological Survey has created ScienceBase, an online 
collaborative scientific data platform. To support collaborative 
data workflows, ScienceBase has four key elements. 1) Data and metadata
cataloging and hosting with options for private, controlled access 
and fully public sharing of data and metadata. 2) Central search and
data discovery including metadata search for linked, non-hosted data. 
3) Full web-service support for all core functionality, including 
standards-based access to specific data types (e.g., Geospatial 
datasets). 4) Research community catalogs to 
enable the organization of data along collaborative group and 
organizational boundaries. 

ScienceBase functionality revolves around the organization of data 
items in a hierarchical structure. These items are flexible representations
of a single dataset with metadata attached. ScienceBase can be used to
store the raw data with attached metadata, or be only the metadata 
repository with links to external data repositories. These data items
are organized into a tree hierarchy (much like the organization of files
on a hard drive). This makes the organization of data by collaborative 
group, project, and/or individual data owner simple. 

Search and access of ScienceBase data items can be accomplished through
a variety of useful interfaces. For manual search and data editing, 
the ScienceBase website offers full functional interfaces. Automated
access is supported by robust REST-ful web services that support all 
functionality offered by the human interfaces. 

The ScienceBase code and infrastructure setup is available to anyone
interested upon request from the ScienceBase team <sciencebase@usgs.gov>.



\begin{itemize}
	\item{What is ScienceBase \url{https://www.sciencebase.gov}}
	\item{Project build for USGS and collaborators}
	\item{Large, flexible, online repository for scientific data}
	\item{Fundamentally hierarchical in design, build on uniquely identified items and children}
	\item{Items have rich metadata functionality that improves data discovery and sharing}
	\item{Build with first-class REST API in mind. All functionality available via web service}
	\item{Available to USGS (DOI?) employees and project collaborators}
\end{itemize}



\section{The \pkg{sbtools} package}
\begin{itemize}
	\item{Goal to allow programmatic access to core SB functionality}
	\item{Harmonize the SB data structures with common R interfaces}
	\item{Lightweight R package that can be quickly added to existing projects and shared broadly}
\end{itemize}


%code example
% \begin{example}
  % x <- 1:10
  % result <- myFunction(x)
% \end{example}

\subsection{Data access API}
The data access functionality of \pkg{sbtools} makes it easy to 
access any public item's data, attached files and metadata. All items
in ScienceBase have a unique identifier that can be used to directly 
access specific items. 

\begin{example}
test_item = item_get("4f4e4b24e4b07f02db6aea14")
test_item
\end{example}

For convenience, \pkg{sbtools} defines an \code{sbitem} object, which is 
returned by \pkg{sbtools} functions when referencing objects. The underlying
data structure is a list. All available metadata for an item can be listed
and accessed in the same way as a named list.

\begin{example}
names(test_item) 
item_get(test_item)$citation

\end{example}

On ScienceBase, all items are organized in a tree structure, with one 
parent and potentially many children. \pkg{sbtools} allows the user to 
easily traverse the tree structure. In some projects, the hierarchy has 
important meaning. For others, it relates only to the owner user or
organization.

\begin{example}

parent = item_get(test_item\$parentId)
parent

item_list_children(parent)

\end{example}

ScienceBase allows items to have attached files. Attached files can easily be 
listed and downloaded directly using \code{item_list_files} and 
\code{item_file_download}. These functions are only useful on items 
that have files attached.


ScienceBase has special support for certain datatypes. 


\begin{itemize}
	\item{Accessing special SB extensions}
	\item{Hierarchy traversal}
\end{itemize}

%code example
% \begin{example}
  % x <- 1:10
  % result <- myFunction(x)
% \end{example}

\subsection{Search API}
\begin{itemize}
	\item{Go over general search API}
	\item{Free text search}
	\item{dateRange query}
	\item{Spatial query}
	\item{Hierarchical SB queries (sub-items)}
	\item{Combination of above}
\end{itemize}

\subsection{ScienceBase authenticated}

\begin{itemize}
	\item{How authentication changes SB}
	\item{Authentication in sbtools}
	\item{Authentication persistence and use throughout package}
	\item{Examples of authentication and other auth-related functions (check and such)}
\end{itemize}

%code example
% \begin{example}
  % x <- 1:10
  % result <- myFunction(x)
% \end{example}

\subsection{Data editing and upload API}
This will be an extensive section
\begin{itemize}
	\item{Item creation}
	\item{Item deletion}
	\item{Attached file upload, replacement and deletion}
	\item{}
\end{itemize}

\subsection{SB Item identifiers}
\begin{itemize}
	\item{Functionality of item identifiers}
	\item{Example of identifier use}
	\item{How identifiers can be useful for projects}
\end{itemize}


\section{Summary}

As many real-world projects have demonstrated, R is
an excellent tool for collaborative and reproducible
research projects \citep{Gandrud2013}. But more flexible and powerful
ways to share and collaborate around common datasets
are a missing component. To help fill this gap, we
have created \pkg{sbtools}. The \pkg{sbtools} package
gives the user programmatic access to the
cloud-based data and metadata storage of USGS ScienceBase
and allows all researchers direct query and download
capabilities to free, public data on ScienceBase.
\pkg{sbtools} enables rapid and reproducible
access to one of the single
largest repositories of earth-science data and an advanced
cloud-based data collaboration platform.



\bibliography{winslow}

\address{Luke A Winslow\\
  U.S. Geological Survey Center for Integrated Data Analytics\\
  Middleton, Wisconsin\\
  USA\\}
\email{lwinslow@usgs.gov}

\address{Scott Chamberlain\\
  rOpenSci / University of California, Berkeley\\
  Berkeley, California\\
  USA\\}
\email{scott@ropensci.org}

\address{Alison Appling\\
  U.S. Geological Survey Center for Integrated Data Analytics\\
  Middleton, Wisconsin\\
  USA\\}
\email{author3@work}

\address{Jordan S Read\\
  U.S. Geological Survey Center for Integrated Data Analytics\\
  Middleton, Wisconsin\\
  USA\\}
\email{jread@usgs.gov}


%some figure code example
%This section may contain a figure such as Figure~\ref{figure:rlogo}.
% \begin{figure}[htbp]
%   \centering
%   \includegraphics{Rlogo}
%   \caption{The logo of R.}
%   \label{figure:rlogo}
% \end{figure}

%code example
% \begin{example}
  % x <- 1:10
  % result <- myFunction(x)
% \end{example}
